\documentclass[10pt,a4paper]{amsart}
\usepackage[margin=19mm]{geometry}
\usepackage{amsmath,amssymb,mathtools}
\usepackage[scr=rsfs,cal=euler]{mathalfa}
\usepackage{xfrac}
\usepackage[unicode,hidelinks,bookmarks=false]{hyperref}
\newcommand{\ie}{i.e.\ }
\newcommand{\cf}{cf.\ }
\newcommand{\resp}{respectively\ }
\newcommand{\Subsection}{\subsection}
\providecommand{\nobreakdash}{\nobreak\hbox{-}}
\setlength{\emergencystretch}{2em}
\multlinegap0pt
\usepackage{tikz}
\newtheorem{lemma}{Lemma}
\newtheorem{theorem}{Theorem}
\title{Bounds for the Vertex Chromatic Number of Connected Triangle-Free Graphs}
\author{S. Akbari$^{a}$, A. Beikmohammadi$^{b}$ \smallskip}
\date{Source: arXiv:2609.07014v1 (CC BY 4.0)}
\begin{document}
\maketitle

\noindent\emph{Source and licence.} Bounds for the Vertex Chromatic Number of Connected Triangle-Free Graphs. Version arXiv:2609.07014v1 (CC BY 4.0), distributed by arXiv under the Creative Commons Attribution 4.0 International licence, http://creativecommons.org/licenses/by/4.0/, as recorded in the arXiv licence metadata for this exact version and shown on the abstract page https://arxiv.org/abs/2609.07014v1. Full licence notice: \texttt{licenses/arxiv-2609.07014-CC-BY-4.0.txt}.\par\smallskip

\noindent\textit{Editorial scope.} Counted unit: the article's theorem ceiling (source0.tex lines 112--118). The article's complete original proof of it is delivered with it (source0.tex lines 120--141, one \texttt{proof} environment of 22 lines). No mathematics is written, repaired or re-derived: the statement and the proof are the article's own lines, in the article's own notation, and no retained line is substituted or re-flowed.  Retained uncounted as the source-local context the proof needs: lines 81--84.  Nothing was omitted from any retained span, so the package carries no omission record.  No retained line contains a citation, so the wrapper carries no bibliography and no bibliography entry.  No retained line contains a figure environment, an image-inclusion command or drawing code, so no image is delivered and the package declares no asset dependency.  Editorial formatting only, in addition: an AMS \texttt{amsart} preamble that restates the article's own declarations and macros used by the retained lines, namely the environments \texttt{lemma}, \texttt{theorem} on separate counters, together with the standard packages those lines need.  The retained environments print in this document as Lemma 1, Theorem 1 in order of appearance, whereas the article prints the same items with the numbering of its own full document.  The complete native source of the article as distributed in the arXiv e-print for this exact version is bundled unchanged as \texttt{sources/209597/source0.tex}.
\begin{lemma} \label{lemma: density}
Let $G$ be a connected $(m,n)$-graph with $\chi(G) = k \geq 3$ and let $H$ be a $k$-critical subgraph of $G$ of order $n_H$. Then
$$m \geq n + \frac{(k-3)\,n_H}{2}.$$
\end{lemma}

\begin{theorem} \label{theorem: no ceiling}
Let $G$ be a connected $(m,n)$-graph of order $n \geq 14$. Then
\[
  \chi(G) \leq \frac{m}{\sqrt{n}} .
\]
Moreover, the bound $14$ is best possible.
\end{theorem}

\begin{proof}
By contradiction suppose that $k = \chi(G) > \frac{m}{\sqrt{n}}$. Then
$$
m < k\sqrt{n} \leq n + \frac{k^2}{4}.
$$
If $k = 1$, then $n = 1$. If $k = 2$, then $n-1 \leq m < 2\sqrt{n}$, so $n \leq 5$.

Let $k \geq 3$ and let $H$ be a $k$-critical subgraph of $G$ of order $n_H \geq k$. By Lemma~\ref{lemma: density} and the fact that $m < n + \frac{k^2}{4}$,
$$
\frac{(k-3)n_H}{2} \leq m-n < \frac{k^2}{4}.
$$
Replacing $n_H$ by $k$ give us $2k(k-3) < k^2$, so $k \in \{3,4,5\}$. We consider these three cases separately, in each case combining the lower bound for $m$ given by Lemma~\ref{lemma: density} with the upper bound $m < k\sqrt{n}$.

If $k=3$, then $n \leq m < 3\sqrt{n}$, so $n \leq 8$.
If $k=4$, then $n_H \geq 4$ gives $m \geq n+2$, so $n + 2 < 4\sqrt{n}$, which implies $n \leq 11$.
If $k=5$, then $n_H \geq 5$ gives $m \geq n+5$, so $n+5 < 5\sqrt{n}$, which implies $n \leq 13$.
In every case $n \leq 13$, a contradiction.

For the sharpness, let $G$ be the graph obtained from $K_5$ by attaching $8$ pendant vertices to one of its vertices. Then $G$ is connected with $n = 13$, $m = 18$ and $\chi(G) = 5$, but
$$\frac{m}{\sqrt{n}} = \frac{18}{\sqrt{13}} \approx 4.993 < 5 = \chi(G).$$
Hence the conclusion of the theorem fails for this graph of order $13$.
\end{proof}

\end{document}
